\documentclass[12pt,aspectratio=169,english,finnish,pdflatex%,handout
]{beamer}
\definecolor{MyGreen}{RGB}{50, 120, 50}
\usecolortheme[named=MyGreen]{structure}

\usepackage{babel}
\usepackage[utf8]{inputenc}
\usepackage[T1]{fontenc}
\usepackage{amsmath,amssymb} 
\usepackage{animate}
\usepackage{multimedia}

\usepackage{listings}
\usepackage{xcolor}
 
\definecolor{codegreen}{rgb}{0,0.6,0}
\definecolor{codegray}{rgb}{0.5,0.5,0.5}
\definecolor{codepurple}{rgb}{0.58,0,0.82}
\definecolor{backcolour}{rgb}{0.95,0.95,0.92}
 
\lstdefinestyle{mystyle}{
    backgroundcolor=\color{backcolour},   
    commentstyle=\color{codegreen},
    keywordstyle=\color{magenta},
    numberstyle=\tiny\color{codegray},
    stringstyle=\color{codepurple},
    basicstyle=\footnotesize,
    breakatwhitespace=false,         
    breaklines=true,                 
    captionpos=b,                    
    keepspaces=true,                 
    numbers=left,                    
    numbersep=5pt,                  
    showspaces=false,                
    showstringspaces=false,
    showtabs=false,                  
    tabsize=2
}
\lstset{style=mystyle}

\usepackage{natbib}
\bibpunct[: ]{(}{)}{,}{}{}{;}

\usepackage{tikz}

\usepackage{tipa}

\usepackage{hyperref}

\setbeamertemplate{navigation symbols}{}

\graphicspath{{figures/}}

\setlength{\leftmargini}{0pt}
\setlength{\leftmarginii}{1em}

%% Write out the names of graphics files being included 
\newwrite\graphics
\immediate\openout\graphics=\jobname.graphics%
\let\oincludegraphics\includegraphics% store original \includegraphics
\renewcommand{\includegraphics}[2][]{% prepend to it (could also use xpatch, etc.)
  \immediate\write\graphics{#2}
  \oincludegraphics[#1]{#2}}

\newcommand\Wider[2][3em]{%
\makebox[\linewidth][c]{%
  \begin{minipage}{\dimexpr\textwidth+#1\relax}
  \raggedright#2
  \end{minipage}%
  }%
}

\newcommand{\kommentti}[1]{
  {\bf[#1]}
}


\title{CSD 598 - Winter 2026
  \\~\\
  An intro to programming and statistical linear model diagnostics
}
\author{Pertti Palo} 
\date{3 March 2026} 


\begin{document}

\frame{\titlepage
  \centering
} 


\frame{\frametitle{Last time: Exploration and inference example}
  Any questions? Any thoughts?
}


\frame{\frametitle{Intro to programming}
  \begin{itemize}
    \item Obviously there is a lot to dig into here.
    \item A sufficient understanding does not require you to know too many details.
    \item Time allowing, we'll look at different programming concepts as we
    work through some ways of diagnosing the behaviour or statistical linear
    models.
    \item First, though, we'll have a look at your questions and possibly use
    some of the following slides to answer them.
    \begin{itemize}
      \item We may also wander into Python or R to try things out.
    \end{itemize}
  \end{itemize}
}


\frame{\frametitle{If you are feeling couragous install these}
    \begin{itemize}
      \item RStudio \url{https://posit.co/downloads/}
      \item Python \url{https://www.python.org/downloads/}
      \item Depending on your operating system some form of app store maybe the
      best bet for installing anything.
      \item If you are super adventurous, install uv
      (\url{https://docs.astral.sh/uv/getting-started/installation/}) and use
      Python that way. (It is actually a really good idea.)
    \end{itemize}
}


\begin{frame}[fragile]
  \frametitle{Basic building blocks I}
  \begin{lstlisting}[language=Python]
  # These are lines comments. They are meant for humans to read
  # and are ignored by the computer.

  # Operators: Do stuff.
  1+1

  # Variables: Make the results stick around.
  a = 1+1
  b = 3
      
  # Combinations thereof: What most programming actually is.
  c = b/a

  # Functions and their siblings: Package recipes for later use.
  square(2)
  \end{lstlisting}
\end{frame}


\frame{\frametitle{Basic building blocks II}
    \begin{itemize}
      \item integer -- aka whole number, 
      \item floating point number -- aka decimal number,
      \item character -- a single character, 
      \item string -- a piece of text made of characters, very useful to us,
      \item lists which will usually contain one given type of variable, e.g.
      a list of numbers or a list of strings.
      \item and many more.
    \end{itemize}
}


\begin{frame}[fragile]
  \frametitle{Basic tools of the trade}
  \begin{itemize}
    \item There are many programming languages and they are suited for
    different tasks and also differ in terms such as readability.
    \item Most relevant to us are R, Python, Matlab, Praat scripts, LibreOffice
    calc/Excel macros.
    \item Code files are usually identified with a suffix that corresponds to
    the language being used: e.g. \verb|.R| for R, \verb|.py| for Python,
    \verb|.praat| for Praat (although the last one is not standardised).
    \item Code can be written with essentially any text editor that will save
    plain text (no formatting frills attached).
  \end{itemize}
\end{frame}


\begin{frame}[fragile]
  \frametitle{Actual tools of the trade}
  \begin{itemize}
    \item Writing code generally easier with an actual code editor or IDE (Integrated
    Development Environment):
    \begin{itemize}
      \item RStudio
      \item Jasp
      \item VSCode (corporate and a bit scary) / VSCodium (probably doesn't
      snitch on you)
      \item PyCharm (often costs money)
      \item Matlab
    \end{itemize}
  \item Consider using an AI helper (like Cursor, Copilot Agent, or Cline), but
  be aware that the less common your task, the more likely they are to produce
  only garbage.
  \item Helpers optimised for actual code writing are probably better than
  general generative models.
  \end{itemize}
\end{frame}


\begin{frame}[fragile]
  \frametitle{Skills of the trade}
  \begin{itemize}
    \item Programming style: Code is read more than written. Learn to write in
    a clear style.
    \item Knowledge of the tools: There will always be new ones. Dedicate some
    time to learning new stuff.
    \item General principles like modularity, 
  \end{itemize}
\end{frame}


\frame{\frametitle{Model diagnosis}
  \begin{itemize}
    \item Collinearity (linear correlation) of the independent/predicting
    variables is bad (mathematically). 
    \begin{itemize}
      \item It can usually be spotted fairly easily by calculating correlations
        among the predictors and/or making scatter plots like a pairs plot.
    \end{itemize}
    \item Heteroscedasticity will mess more the efficiency of prediction than
    the accuracy of it.
    \item Outliers can mess prediction accuracy.
    \item Here a bit of data exploration and model diagnostic plots go a long
    way:
    \begin{itemize}
      \item Neat explanation of diagnostic plots can be found here
      \url{https://library.virginia.edu/data/articles/diagnostic-plots}
      \item Collinearity can be seen in for example a pairs plot or a
      correlation matrix (not to be confused with e.g. "Correlation of Fixed
      Effects" in `lmer' output).
    \end{itemize}
  \end{itemize}
}


\frame{\frametitle{Some resources}
  \begin{itemize}
    \item RStudio \url{https://posit.co/downloads/}
    \item Jasp \url{https://jasp-stats.org/}
    \item Python learning materials
    \url{https://wiki.python.org/moin/BeginnersGuide/NonProgrammers}
    \begin{itemize}
      \item Especially the Interactive Courses are really good for learning programming in general.
    \end{itemize}
    \item Praat scripting guide \url{https://praatscriptingtutorial.com/}
  \end{itemize}
}


% \frame{\frametitle{References}
  
% \bibliographystyle{apalike}
% \bibliography{science_combined}

% }
\end{document}

